% !TeX root = ../Thesis.tex

%*******************************************************
% Abstract
%*******************************************************

\iftoggle{cts@phd}{%
\begingroup
\let\clearpage\relax
\let\cleardoublepage\relax
\let\cleardoublepage\relax
}{}

\addchapCondToC{Abstract}
Diese Arbeit untersucht, ob sich einzelne visuelle Merkmale auf antiken Münzbildern ohne münzspezifisches Training detektieren und für den Vergleich von Münzen nutzen lassen. Manuelle Stempelstudien sind sehr aufwendig, und bestehende bildbasierte Verfahren erfordern meist ein Training auf münzspezifischen Daten.
Für die Detektion wird das Open-Vocabulary-Modell YOLO-World eingesetzt, das Merkmale im Zero-Shot-Verfahren anhand von Textbeschreibungen erkennt. Zu jedem detektierten Merkmal wird ein Feature-Vektor aus einer internen Schicht des Netzwerks extrahiert. Merkmale derselben Klasse auf verschiedenen Münzen werden mit euklidischer Distanz, Manhattan-Distanz und Kosinus-Ähnlichkeit verglichen. Die Umsetzung erfolgt als Anwendung, die Verarbeitung, Speicherung und Suche ermöglicht.
Die Auswertung nutzt 200 Münzen aus dem Corpus Nummorum in 20 Stempelgruppen zu je zehn Münzen als Grundwahrheit. Geprüft wird, ob Münzen mit demselben Vorderseitenstempel unter den ähnlichsten Treffern liegen (HitRate@1, HitRate@5, Precision@5). YOLO-World liefert bei einem Schwellenwert von 0{,}005 auf den Vorderseiten 905 Detektionen mit meist niedriger Confidence, nur 23 der 57 Prompts treten als Klasse auf. Über alle Klassen liegt bei 30,8\,\% der Anfragen eine Münze desselben Stempels unter den fünf ähnlichsten Treffern, bei zufälliger Reihenfolge wären es 21,2\,\%. Deutlich über dem Zufall liegen nur Gesichter (\textit{human face}: 42,0\,\% gegenüber 20,9\,\%), bei anderen Klassen trennen die Vektoren nicht. Da auf Münzen desselben Stempels oft verschiedene Klassen detektiert werden, bleiben viele stempelgleiche Münzen unverglichen. Das ohne Open Vocabulary arbeitende YOLOv8n erreicht vergleichbare Werte. Grenzen ergeben sich aus der Abhängigkeit von der Prompt-Liste, der kleinen Datenbasis und dem Erhaltungszustand der Münzen.

\iftoggle{cts@phd}{%
\vfill
\begin{otherlanguage}{ngerman}
  \addchapCondToC{Zusammenfassung}
  Kurze Zusammenfassung des Inhaltes in deutscher Sprache\dots
  \lipsum[11-20][1-12]
\end{otherlanguage}
\endgroup
\vfill
}{}
